\section{Introduction}
\label{sec:intro}

The parameter scale of modern transformer architectures has outpaced the physical memory capacities of commodity edge and client computing devices. While state-of-the-art dense and sparse Mixture-of-Experts (MoE) models offer unprecedented reasoning capabilities, executing a 7-billion parameter dense model or an MoE architecture with dozens of routed experts requires between 4\,GB and 32\,GB of high-speed memory. On commodity laptops and embedded systems—typically equipped with 4\,GB to 8\,GB of unified system RAM and commodity SATA or entry-level NVMe storage (delivering 300--500\,MB/s sequential bandwidth)—running these models natively has historically been deemed intractable.

Existing systems address this memory deficit through operating system virtual memory swapping or memory-mapped file access (\texttt{mmap})~\citep{sheng2023flexgen,gerganov2023llama}. However, traditional OS-level demand paging suffers from structural pathologies when applied to autoregressive inference:
\begin{enumerate}
    \item \textbf{Random Access Amplification}: In MoE architectures, token-level routing dynamically selects a sparse subset of expert feed-forward networks (e.g., top-2 of 16, or top-8 of 64) at every layer. Because standard weight formats store tensors sequentially by layer and expert index, accessing routed experts causes scattered 4\,KB random reads across storage blocks, collapsing commodity SSD throughput down to single-digit megabytes per second.
    \item \textbf{Page Fault Stalls and Thrashing}: Standard \texttt{mmap} implementations rely on the OS kernel's page cache and virtual memory management subsystem. When available physical RAM is significantly smaller than the active working set, page replacement algorithms (such as global active/inactive list LRU) evict pages that are immediately required by subsequent tokens, triggering severe page-fault thrashing, kernel lock contention on \texttt{mm->mmap\_lock}, and eventual invocation of the Linux Out-Of-Memory (OOM) killer.
    \item \textbf{Kernel Double-Buffering Overhead}: Standard cached POSIX reads route storage blocks through kernel page-cache buffers before copying data into user-space process memory, polluting CPU L2/L3 caches and adding redundant memory copy operations on memory-bandwidth-starved architectures.
\end{enumerate}

To overcome these barriers, we formulate the \textbf{Commodity SATA Thesis}: \emph{By co-designing physical on-disk weight layouts with user-space asynchronous buffered I/O pipelines and empirical residency heuristics, large language models can be reliably streamed and executed on low-end commodity client hardware under strict OS memory limits that trigger fatal OOM failures in standard virtual memory systems.}

In this paper, we present \textbf{Kura}, a storage-native decoupled execution engine designed specifically for severe memory constraints. Kura makes three central systems contributions:
\begin{itemize}
    \item \textbf{Physical Layout Coalescing (\textsc{Loom})}: We design Loom, an offline layout compiler that measures expert co-activation frequencies across calibration sequences, constructs a sparse Markov graph, and solves a 1D graph arrangement problem via continuous spectral relaxation. Loom physically serializes frequently co-activated experts into contiguous storage extents, converting scattered random seeks into broad sequential reads.
    \item \textbf{User-Space Buffered I/O Arenas (\textsc{RamArena})}: Rather than relying on OS page faults or unaligned kernel-bypass direct DMA, Kura pre-allocates fixed, 64-byte aligned user-space memory slabs and orchestrates non-blocking prefetching using Linux \texttt{io\_uring} and POSIX \texttt{pread64}. This eliminates runtime page allocations in the execution hot path and ensures predictable resident set size (RSS) stability.
    \item \textbf{Progressive Tiered-Fidelity Previews (\textsc{Ember})}: We evaluate Ember, an execution scheme based on offline precompiled truncated Singular Value Decomposition (SVD) expert sketches. In precompiled Mode C, Ember synthesizes a provisional speculative draft preview token in $0.16$\,ms (exhibiting $1.99\%$ relative $L_2$ divergence) to provide immediate user-facing responsiveness while exact weights load asynchronously over $739.40$\,ms.
\end{itemize}

Crucially, this work also documents the \textbf{Economics of Systems Overhead}: through adversarial ablation, we show that complex online dynamic optimization frameworks—specifically KAEF's convex shadow-pricing auction scheduler—incur an average controller overhead of $\approx 12.35$\,ms per decision ($4,741.47$\,ms cumulative sequence latency). At intermediate memory capacities, this optimization overhead exceeds the storage latency saved by prefetching, making static LRU combined with Loom layout compilation the superior production choice. Similarly, we demonstrate that attempting to compute SVD low-rank sketches dynamically online introduces a catastrophic $3,018.20$\,ms CPU bottleneck ($+408.2\%$ latency regression), proving that progressive fidelity mechanisms must remain strictly decoupled from the live serving path.
